\begin{frame}{Conclusão e trabalhos futuros}
  \begin{itemize}
    \item[$\blacktriangleright$] Conclusão
      \begin{itemize}
        \item As técnicas de re-amostragem mostraram-se eficazes para lidar com o desbalanceamento;
        \item Abordagens custo-sensíveis dependem fortemente da definição adequada dos custos de classificação;
        \item Não existe uma solução única para todos os contextos de crédito, e a escolha da técnica deve ser guiada pelas características específicas do conjunto de dados, como o grau de desbalanceamento, o número de amostras e a dimensionalidade dos dados.
      \end{itemize}
  \end{itemize}
\end{frame}

\begin{frame}{Conclusão e trabalhos futuros}
  \begin{itemize}
    \item[$\blacktriangleright$] Trabalhos futuros
      \begin{itemize}
        \item Explorar outras técnicas de re-amostragem, como o ADASYN e o SMOTE-ENN;
        \item Investigar a eficácia de técnicas de meta-aprendizado, como o MetaCost;
        \item Deep-Learning.
      \end{itemize}
  \end{itemize}
\end{frame}
